\section{Retrieval error and the transmission profile}\label{sec:results-transfer}

The transfer theorem makes the ratio of component error to transmission the natural retrieval
diagnostic, and the conditional-profile theorem asks for the profile $t\mapsto\E[e_R^2\mid t]$.
This section measures both on finite test sets; a slope fitted to a finite profile is a
description of that profile, not an asymptotic rate.

Table~\ref{tab:v2-transmission} uses the three-member width-2000 pipeline on splits 104--107,
with float64 predictions of nine heads on 2{,}331 test states and 285 bands, or 664{,}335
state--band entries per split. The retrieval uses $\rho=0.7$, $R=0.9$ and an albedo bound $S$
taken from the training data.

\begin{table}[p]
\centering
% Generated by code/format_publication_tables.py; numerical cells preserved.
\textbf{A. Errors and finite-range profiles}\par\medskip
{\fontsize{9}{11}\selectfont
\setlength{\tabcolsep}{3pt}
\begin{tabular}{lrrrr}
\toprule
head & $\varepsilon_R$ & $2\nu$ (fit) & \thead{raw p95\\raw p99} & \thead{constrained p95\\constrained p99} \\
\midrule
isotropic kernel & 0.0886$\pm$0.0331 & 0.16$\pm$0.11 & \shortstack{0.696$\pm$0.000\\0.97$\pm$0.02} & \shortstack{0.266$\pm$0.018\\0.70$\pm$0.00} \\
network & 0.1650$\pm$0.0177 & 0.20$\pm$0.16 & \shortstack{0.670$\pm$0.026\\0.99$\pm$0.14} & \shortstack{0.268$\pm$0.111\\0.70$\pm$0.00} \\
network ensemble & 0.1085$\pm$0.0254 & 0.18$\pm$0.14 & \shortstack{0.625$\pm$0.013\\0.86$\pm$0.08} & \shortstack{0.239$\pm$0.078\\0.70$\pm$0.01} \\
network + residual & 0.0698$\pm$0.0373 & 0.21$\pm$0.14 & \shortstack{0.489$\pm$0.054\\0.89$\pm$0.01} & \shortstack{0.200$\pm$0.000\\0.70$\pm$0.00} \\
ensemble + residual & 0.0690$\pm$0.0381 & 0.22$\pm$0.14 & \shortstack{0.464$\pm$0.033\\0.88$\pm$0.01} & \shortstack{0.200$\pm$0.000\\0.70$\pm$0.00} \\
feature kernel & 0.0514$\pm$0.0454 & 0.26$\pm$0.16 & \shortstack{0.075$\pm$0.014\\0.89$\pm$0.05} & \shortstack{0.076$\pm$0.015\\0.70$\pm$0.00} \\
concatenated features & 0.0634$\pm$0.0680 & 0.28$\pm$0.16 & \shortstack{0.060$\pm$0.004\\0.92$\pm$0.03} & \shortstack{0.061$\pm$0.004\\0.70$\pm$0.00} \\
coordinate selection & 0.0721$\pm$0.0648 & 0.25$\pm$0.15 & \shortstack{0.098$\pm$0.031\\0.85$\pm$0.03} & \shortstack{0.099$\pm$0.032\\0.70$\pm$0.00} \\
convex stack & 0.0600$\pm$0.0661 & 0.25$\pm$0.14 & \shortstack{0.143$\pm$0.015\\1.09$\pm$0.05} & \shortstack{0.151$\pm$0.015\\0.70$\pm$0.00} \\
\bottomrule
\end{tabular}}
\par\bigskip\textbf{B. Failure counts and legacy diagnostics}\par\medskip
{\fontsize{9}{11}\selectfont
\setlength{\tabcolsep}{3pt}
\begin{tabular}{lrrrr}
\toprule
head & \thead{failure [\%]\\upper clip [\%]} & \thead{$\E\min\{R^2,$\\$e_R^2/t^2\}$} & \thead{legacy\\violations} & \thead{diagnostic\\bound / MSE} \\
\midrule
isotropic kernel & \shortstack{1.09$\pm$0.03\\3.56$\pm$0.07} & 0.064$\pm$0.000 & 0 & 3.9$\pm$0.2 \\
network & \shortstack{1.31$\pm$1.00\\2.93$\pm$0.43} & 0.056$\pm$0.003 & 0 & 3.4$\pm$0.7 \\
network ensemble & \shortstack{0.61$\pm$0.33\\2.56$\pm$0.92} & 0.051$\pm$0.002 & 0 & 3.7$\pm$0.9 \\
network + residual & \shortstack{0.82$\pm$0.02\\2.51$\pm$0.12} & 0.044$\pm$0.001 & 0 & 4.0$\pm$1.2 \\
ensemble + residual & \shortstack{0.85$\pm$0.01\\2.43$\pm$0.07} & 0.043$\pm$0.001 & 0 & 3.9$\pm$1.3 \\
feature kernel & \shortstack{0.71$\pm$0.02\\1.53$\pm$0.10} & 0.031$\pm$0.001 & 0 & 5.1$\pm$2.3 \\
concatenated features & \shortstack{0.72$\pm$0.02\\1.49$\pm$0.02} & 0.030$\pm$0.000 & 0 & 5.3$\pm$2.4 \\
coordinate selection & \shortstack{0.73$\pm$0.02\\1.61$\pm$0.13} & 0.033$\pm$0.002 & 0 & 4.7$\pm$2.0 \\
convex stack & \shortstack{0.84$\pm$0.05\\1.78$\pm$0.08} & 0.036$\pm$0.001 & 0 & 4.4$\pm$2.0 \\
\bottomrule
\end{tabular}}

\caption{Transmission diagnostics for the three-member width-2000 pipeline, splits 104--107.
Each cell with a spread is a mean $\pm$ sample standard deviation over the four splits. All
reflectance errors in this table are in reflectance units, not percentage points. Panel A gives
the retrieval-weighted error, the fitted profile slope and the raw and constrained percentiles;
panel B gives failure and clipping rates and the bound diagnostics. The violation count tests
the pointwise bound on the entries with $s\le S$; on entries with nonpositive transmission or an
albedo outside $[0,1)$ the bound expressions are diagnostics rather than bounds.}
\label{tab:v2-transmission}
\end{table}

\paragraph{The physical domain in these splits.}
$S$ is the largest finite training albedo not exceeding one, and the training blocks hold 227 or
454 entries above one depending on the split. On split 104, 228 of the 664{,}335 test entries
exceed $S$, and none do on the other three splits; there are 203, 6, 3 and 3 nonpositive
transmissions on splits 104, 105, 106 and 107. The pointwise check covers the entries with
$s\le S$, and the expectations and ratios in the table are averages over all entries.
Section~\ref{sec:domain} scores the main pipeline on the physical domain itself.

\paragraph{Observed error profiles.}
The feature kernel has $\varepsilon_R=0.0514$ and the single network $0.1650$. The mean fitted
profile slopes $2\nu$ range from $0.16$ to $0.28$. The raw 95th-percentile error is $0.075$ for
the single-member feature kernel and $0.060$ for concatenated features, against $0.143$ for the
stack, and the constrained 95th percentiles are $0.076$, $0.061$ and $0.151$. The ratio of the
middle bound to the mean squared error lies between about $3.4$ and $5.3$ across heads.

The middle term of Theorem~\ref{thm:transfer}(ii), $\E\min\{R^2,e_R^2/t^2\}$, orders the nine
heads by their retrieval tails. Taking the four-split means, its rank correlation with the raw
95th percentile is exactly one (Spearman and Kendall), and with the constrained 95th percentile
$0.98$ (Spearman), while the unweighted $\varepsilon_R$ reaches only $0.78$ and $0.81$. The stack,
for instance, has a smaller $\varepsilon_R$ than the kernel on concatenated features ($0.060$
against $0.063$) but a larger middle term ($0.036$ against $0.030$) and a constrained tail more
than twice as large. What orders the retrievals is the component error weighted by the
transmission, as the theorem says, not the component error itself.

About five percent of the entries have transmission below $3\times10^{-4}$ of the training flux
scale, and among positive transmissions the one-percent quantile is of order $10^{-10}$ of that
scale, while the share of squared retrieval-weighted component error on entries below $0.03$ of
the scale is under a thousandth for every head. Small absolute errors are amplified by weak
transmission. The heads still differ in the inverse tail, because the location, sign and
dependence of their errors matter, as the exact identities show.

The table's own transmission distribution is extremely heavy near zero: relative to the median
positive flux, its $1\%$, $3\%$, $5\%$ and $10\%$ quantiles are $1.1\times10^{-10}$,
$3.6\times10^{-6}$, $2.2\times10^{-4}$ and $1.7\times10^{-2}$, a log-log slope of
$\beta\approx0.12$--$0.13$ for $F_t$ over the lowest tenth of the entries. With the measured
profile slopes $2\nu\approx0.16$--$0.28$, Theorem~\ref{thm:profile} places the all-band retrieval in
its slow regime, $\beta<2(1-\nu)$, where the mean squared error scales like
$\sigma^{\beta/(1-\nu)}$ with exponent $0.13$--$0.15$: halving the component error would shrink it
by less than a tenth. On a passband $t\ge t_0>0$ fixed in advance the bound of item (iv) is
linear in $\varepsilon_R$ instead. A slope fitted over a finite range of transmissions does not
establish the tail law down to zero or an envelope uniform over training sizes, so this reading
describes the table's regime rather than proving a rate. Scaling a fixed error profile down buys
little over all bands; moving the error away from the opaque entries, as the feature kernels do,
buys much more.

\section{The stacking objective on data}\label{sec:results-stacks}

Table~\ref{tab:v2-stacks} compares stacks fitted under different objectives on five splits:
the wide pipeline at seed 101 and the three-member pipeline on splits 104--107. The objectives
and the projection sequence are those of Section~\ref{sec:joint}.

\begin{table}[p]
\centering
\input{layout_v2_stacks}
\caption{Stacks under different fitting objectives at $\rho=0.7$, five splits. The
finite-threshold rows use the separable quadratic $J'_\tau$ of \eqref{eq:Jtau-sep}, with
$u=\tau/T_{\rm train}$; the unweighted row is a component squared-error baseline. Radiance is in
percent and the reflectance percentiles in reflectance units. Weights are fitted approximately,
as described in Section~\ref{sec:joint}.}
\label{tab:v2-stacks}
\end{table}

The mean-relative-norm stack has radiance error $0.0754\%$ and constrained 95th-percentile
error $0.141$, and the row-relative squared-error stack $0.1130\%$ and $0.180$. The
finite-threshold $J'_\tau$ fits give radiance errors of about $0.22$--$0.23\%$ for $u\le0.1$,
$0.1491\%$ at $u=0.3$ and $0.1175\%$ at $u=1$, and the unweighted baseline $0.1262\%$.
Weighting by the transmission moves the stack the way Proposition~\ref{prop:stackweights}
predicts: at $u=0.003$ the constrained 95th percentile falls from $0.141$ to $0.103$ while the
radiance error rises from $0.075\%$ to $0.216\%$. It does not reach the kernel on features alone
($0.077$) or on concatenated features ($0.061$). The fitting loss and the reported metric need
not favor the same combination, and $J'_\tau$ is a mean-square criterion while the table reports
a quantile; the factor-four transfer bound holds for every feasible weight vector on a
distribution satisfying its hypotheses.

